\begin{ledgroupsized}[r]{120mm}
\footnotesize
\pstart
\noindent\textbf{Überlieferung:}
\pend
\end{ledgroupsized}
\begin{ledgroupsized}[r]{114mm}
\footnotesize
\pstart
\parindent -6mm
\makebox[6mm][l]{\textit{L}}%
Konzept:
LH~XXXV 10, 16 Bl.~1\textendash2.
Ein Bogen~2\textsuperscript{o};
\href{\Scan{LH_35_10_16_001-002wz}}%
{Wasserzeichen} im Falz.
Vier Seiten.
\pend
\end{ledgroupsized}
\selectlanguage{latin}
\vspace{8mm}
\pstart
\noindent
\href{\Scan{LH_35_10_16_001r}}%
{\lbrack1~r\textsuperscript{o}\rbrack}\  
\hfill
\edlabel{35_10_16_001-002_6a}%
Viennae\indexO{Wien} mense Februario 1713
\pend
\pstart
\noindent
(1) Cum nihil hactenus scriptorum celeberrimi viri 
Hieronymi Saccherii\indexN{\textls{Saccheri}, Girolamo 1667\textendash1733} 
\edtext{vidissem, eo magis 
\edtext{\cite{00541}\textit{Neostatica} ejus jam Mediolani\indexO{Mailand} A.\,D.~1708 edita,}{\lemma{\textit{Neostatica} \lbrack...\rbrack\ edita}\Cfootnote{\indexN{\textls{Saccheri}, Girolamo 1667\textendash 1733}\textsc{G.~Saccheri}, \cite{00541}\textit{Neo-statica}, Mailand 1708.}}
sed nunc demum a R\textsuperscript{mo} Patre Granello\indexN{\textls{Granelli} (Granellus), Carlo 1671\textendash1739} 
mihi communicata\edlabel{35_10_16_001-002_6b}		%
sum delectatus in qua}{\lemma{vidissem,}\Bfootnote{%
\textit{(1)}~gavisus sum %
\textit{(2)}~eo magis \textit{Neostatica} ejus
\bars{jam Mediolani A.\,D.~1708 edita, sed nunc demum %
\textit{(1)}~mihi %
\textit{(a)}~visa, %
\textit{(b)}~cognita %
\textit{(2)}~a R\textsuperscript{mo} Patre Granello mihi communicata~\textit{erg.}}
sum delectatus %
\textit{(a)}~plena ingenii, et %
\textit{(b)}~in qua~\textit{L}}} 
multa sunt ingenii acris, meditationisque profundae specimina. Probat 
\edtext{ut video}{\lemma{}\Bfootnote{ut video \textit{erg.~L}}} 
\edtext{calletque Geometriam infinitesimalium,\indexS{Geometria infinitesimalium} quae}{\lemma{calletque}\Bfootnote{%
\textit{(1)}~egregie calculum
\textit{(a)}~infinitesi %
\textit{(b)}~infinitesimalium %
\textit{(2)}~Geometriam infinitesimalium \bars{(de qua \textit{Neostaticae} pag.~79) \textit{erg.~u.~gestr.}}, quae~\textit{L}}}
\edtext{calculo\indexS{calculus infinitesimalium} quodam introducto a me nonnihil fuit promota,}{\lemma{calculo \lbrack...\rbrack\ promota}\Cfootnote{\textsc{G.\,W.~Leibniz}, \glqq Nova methodus pro maximis et minimis\grqq, \textit{AE} Oktober 1684, S.\;467\textendash473.}}
\edtext{ut ipse agnoscit \textit{Neostaticae} pagina 79.}{%
\lemma{ut \lbrack...\rbrack\ 79}%
\Cfootnote{%
\cite{00541}a.\,a.\,O., lib.~3, Praenotanda, Nr.~5, S.\;79. Dort verweist \indexN{\textls{Saccheri}, Girolamo 1667\textendash 1733}Saccheri außerdem auf einige \glqq clarissimos Geometras\grqq, allen voran \glqq Liebnitium\grqq\ (\textit{sic!}).}}
\pend
\begin{edfigure}%
\hspace*{4em}%
\includegraphics[scale=1]{%
gesamttex/edit_VIII-4/images/LH_35_10_16_001-002_d1_001r.pdf}
\edcaption{\lbrack\textit{Fig.~1}\rbrack}
\end{edfigure}
\pstart
Molitur etiam non vulgaria ad elementa\indexS{elementa rei staticae}\indexS{elementa compositionis motuum} quaedam rei staticae\indexS{res staticae} et compositionis motuum\indexS{compositio motuum} demonstranda. 
\edlabel{35_10_16_001-002_7a}%
Qua 
\edtext{occasione commemorare}{\lemma{occasione}\Bfootnote{%
\textit{(1)}~memorabo %
\textit{(2)}~memorare %
\textit{(3)}~commemoro %
\textit{(4)}~commemorare~\textit{L}}} 
placet, 
\edtext{olim a me in Parisino \textit{Eruditorum diario}
\textit{Generalem}\indexS{regula generalis de compositione motuum}  quandam \textit{regulam compositionis}
\edtext{\textit{motuum} fuisse}{\lemma{\textit{motuum}}\Bfootnote{%
\textit{(1)}~regulam %
\textit{(2)}~fuisse~\textit{L}}}
editam,}{\lemma{olim \lbrack...\rbrack\ editam}%
\applabel{FnC:35_10_16_001-002_regle}%
\Cfootnote{\indexN{\textls{Leibniz} (Leibnitius, GGL), Gottfried Wilhelm 1646\textendash1716}\textsc{G.\,W.~Leibniz},
\cite{00542}\textit{Regle generale de la composition des mouvemens}, \cite{00157}\textit{JS} vom 7.~September 1693, S.\;417\textendash419.}}
quae huc redit. %
\pend
\pstart
Mobile\indexS{mobile} \textit{M} 
\edtext{puncti instar consideratum\indexS{mobile instar puncti consideratum} 
(nempe sumto in mobilis locum ipso ejus centro gravitatis\indexS{centrum gravitatis})}{\lemma{}\Bfootnote{puncti instar consideratum (nempe sumto %
\textit{(1)}~ejus loco %
\textit{(2)}~in mobilis locum ipso ejus centro gravitatis) \textit{erg.~L}}} 
feratur motu composito\indexS{motus ex motibus componentibus quotcunque compositus} ex motibus componentibus\indexS{motus componens} quotcunque,
\edtext{velut}{\lemma{}\Bfootnote{velut \textit{erg.~L}}} 
\textit{MF}, \textit{MG}, 
\edtext{\textit{MH}\lbrack,\rbrack\ quarum rectarum unaquaque repraesentari ponitur directio et velocitas\indexS{directio et velocitas motus} cujusque motus. Fingatur mobile ferri}{\lemma{\textit{MH}}\Bfootnote{%
\bars{%
\textit{(1)}~quibus rectis repraesentari ponantur directiones et velocitates %
\textit{(2)}~quarum \lbrack...\rbrack\  motus. \textit{erg.}}
Fingatur
\bars{mobile \textit{erg.}}
ferri~\textit{L}}} 
unoquoque separatim,\indexS{mobile latum unoquoque motu separatim} tanquam esset multiplicatum, et simul veniret 
\edtext{in \textit{F}, \textit{G}, \textit{H}. Quaeratur}{\lemma{in \textit{F}, \textit{G}, \textit{H}. }\Bfootnote{%
\textit{(1)}~Quaeratur centrum \textit{K} %
\textit{(2)}~Quaeratur~\textit{L}}} 
\textit{K}, centrum
\edtext{commune\indexS{centrum commune} omnium horum}{\lemma{commune}\Bfootnote{%
\textit{(1)}~horum trium %
\textit{(2)}~omnium horum~\textit{L}}} 
punctorum seu 
\edtext{locorum quae Mobile}{\lemma{locorum}\Bfootnote{%
\textit{(1)}~quae %
\textit{(2)}~quae Mobile~\textit{L}}} 
unoquoque motu componente\indexS{motus componens} intra idem 
\edtext{tempus nancisceretur:}{\lemma{tempus}\Bfootnote{%
\textit{(1)}~mobile 
\textit{(2)}~nas %
\textit{(3)}~nancisceretur:~\textit{L}}} 
his positis feretur mobile
directione ipsius \textit{MK}, viae centri 
\edtext{communis,\indexS{via centri communis} sed directione simul et}{\lemma{communis,}\Bfootnote{%
\textit{(1)}~et %
\textit{(2)}~\bars{sed \textit{erg.}}  directione simul %
\textit{(a)}~ac %
\textit{(b)}~et~\textit{L}}} 
velocitate ut \textit{ML}, quae sit ad \textit{MK}, ut numerus motuum componentium\indexS{numerus motuum componentium} ad unitatem.%
\edlabel{35_10_16_001-002_7b}%
\pend 
\pstart 
(2) Vellem expendere totum 
\edtext{opus per praesentes distractiones liceret, haud dubie enim non mediocrem inde fructum percepturus essem. 
\edlabel{35_10_16_001-002_8a}%
Cursoria tamen}{\lemma{opus}\Bfootnote{%
\textit{(1)}~liceret per praesentes distractiones, cursoria tamen %
\textit{(2)}~liceret, \lbrack...\rbrack\ tamen~\textit{L}}}
lectione alicubi haesi. Nempe 
\edtext{video 
\edtext{asseverari in corollario propositionis 15, lib.~\lbrack2.\rbrack, solicitationes}{%
\lemma{asseverari}%
\Bfootnote{%
\textit{(1)}~lib.~1. prop.~15. corollario, %
\textit{(2)}~co
\textit{(3)}~in corollario propositionis 15, lib.~\bars{1.~\textit{ändert Hrsg.}}, solicitationes~\textit{L}}}
gravium\indexS{solicitatio gravis} ex quiete versus centrum\indexS{solicitationes gravium versus centrum} gravium commune\indexS{centrum gravium commune}
non posse esse in quavis distantia aequales\indexS{solicitationes in quavis distantia aequales} et multominus in minore distantia posse 
\edtext{esse}{\lemma{}\Bfootnote{esse \textit{erg.~L}}} 
minores, sed potius in majore distantia debere esse majores.}{\lemma{video \lbrack...\rbrack\ majores}%
\Cfootnote{\indexN{\textls{Saccheri}, Girolamo 1667\textendash 1733}%
\textsc{Saccheri}, \cite{00541}\textit{Neo-statica}, lib.~2, prop.~15, \textit{Corollarium}, S.\;69.}}
Fateri autem cogor mihi contrarium persuasum fuisse; usque adeo ut olim
\edtext{in \textit{Tentamine de Motuum coelestium causis} posuerim solicitationes paracentricas\indexS{solicitatio paracentrica} esse reciproce ut quadrata distantiarum;\indexS{quadratum distantiae}\indexS{solicitationes paracentricae reciprocae quadratis distantiarum} eo fere
\href{\Scan{LH_35_10_16_001v}}%
{\lbrack1~v\textsuperscript{o}\rbrack}\  
argumento quo statuunt dudum Optici,\indexS{Optici} illuminationes\indexS{illuminatio}\indexS{illuminationes reciprocae quadratis distantiarum ab illuminante} esse reciproce ut quadrata distantiarum\indexS{quadratum distantiae} ab illuminante:\indexS{illuminans} 
\edlabel{35_10_16_001-002_8b}%
dum radios solicitantes\indexS{radius solicitans} perinde concipio 
\edtext{ac illuminantes.\indexS{radius illuminans}}{%
\lemma{ac}%
\Bfootnote{%
\textit{(1)}~illuminandos.
\textit{(2)}~illuminantes.~\textit{L}}}%
}{\lemma{in \textit{Tentamine} \lbrack...\rbrack\ illuminantes}%
\Cfootnote{
Siehe die unveröffentlichte sog. \glqq zweite Bearbeitung\grqq\
des \cite{00543}\textit{Tentamen de motuum coelestium causis}, bspw. LH~XXXV~9, 2~Bl.~54~v\textsuperscript{o} und \cite{01043}\textit{LMG}~VI, S.\;165\,f.
Vgl. die Anmerkung zu N.\;\ref{60124}, \refnote{FnC:35_10_16_004_Tentamen}.}} 
Ut rem a physica consideratione ad mathematicam revocem, seponamque modum per quem solicitatio paracentrica\indexS{solicitatio paracentrica} procuratur. 
\edtext{Caeterum considerationem qua centrum gravium\indexS{centrum gravium} sumimus ut finite distans,\indexS{distantia finita} seu tendentias\indexS{tendentiae convergentes} ut convergentes, 
Neostaticam\indexS{consideratio Neostatica} cum cl. Autore appellare licebit.}{\lemma{}\Bfootnote{Caeterum considerationem qua centrum gravium %
\textit{(1)}~consideramus %
\textit{(2)}~sumimus \lbrack...\rbrack\  Neostaticam \bars{cum cl. Autore \textit{erg.}} appellare licebit. \textit{erg.~L}}} 
\pend
\begin{edfigure}%
\includegraphics[scale=1]{%
gesamttex/edit_VIII-4/images/LH_35_10_16_001-002_d2_001r.pdf}
\edcaption{%
\edtext{\lbrack\textit{Fig.~2}\rbrack}{\lemma{%
\lbrack\textit{Fig.~2}\rbrack}\Cfootnote[nonum]{Tlw. Übernahme der Vorlage: 
\indexN{\textls{Saccheri}, Girolamo 1667\textendash 1733}\textsc{Saccheri}, \cite{00541}\textit{Neo-statica}, lib.~2, Postulatum, S.\;56. Zwei gestrichene Entwürfe zum Diagramm werden nicht wiedergegeben.}}}
\end{edfigure}
\pstart 
(3) Rem autem possibilem esse, ut solicitationes\indexS{solicitatio}\indexS{solicitatio aucta distantia decrescens} decrescant aucta distantia a centro,\indexS{distantia a centro} sic videtur ostendi 
\edtext{posse: sit recta \textit{RC} inque ejus extremis \textit{R} et \textit{C} sint frusta ferrea;\indexS{frustum ferreum} in \textit{D} autem sit magnes.\indexS{magnes} Constat}{\lemma{posse:}\Bfootnote{%
\textit{(1)}~sint %
\textit{(a)}~in %
\textit{(b)}~\textit{R} et 
\textit{(aa)}~\textit{C}
\textit{(bb)}~in
\textit{(cc)}~\textit{C}
\textit{(2)}~sit recta \textit{RC} inque ejus extremis \textit{R} et %
\bars{sint \textit{erg.}}
frusta ferrea; in \textit{D} \bars{autem sit \textit{erg.}} magnes. %
\textit{(aa)}~Manifestum
\textit{(bb)}~Constat~\textit{L}}} 
magnetem\indexS{magnes} minus trahere ex distantia majore. Et 
\edtext{tamen poterunt}{\lemma{tamen}\Bfootnote{%
\textit{(1)}~poterit %
\textit{(2)}~posito angulo \textit{DNR} %
\textit{(a)}~\bars{recto \textit{erg.}}, poterunt %
\textit{(b)}~recto, %
\textit{(3)}~poterunt~\textit{L}}} 
ita 
\edtext{assignari frustorum ferreorum\indexS{frustum ferreum} et brachiorum\indexS{brachium} \textit{NC}, \textit{NR} magnitudines, ut \textit{N} sit}{\lemma{assignari}\Bfootnote{%
\textit{(1)}~rationes %
\textit{(2)}~frustorum ferreorum %
\textit{(a)}~et bra %
\textit{(b)}~inter se, %
\textit{(c)}~et brachiorum \textit{NC}, \textit{NR}, %
\textit{(aa)}~ut \textit{N} sit %
\textit{(bb)}~magnitudines, ut \textit{N} sit~\textit{L}}} 
punctum aequilibrii.\indexS{punctum aequilibrii} Possunt autem 
\edtext{frus\lbrack ta\rbrack}{%
\lemma{}%
\Bfootnote{%
frustra 
\textit{L ändert Hrsg.}}}
ferrea\indexS{frustum ferreum} conferri ponderibus, magnes\indexS{magnes} \textit{D} centro gravium.\indexS{centrum gravium} 
\pend
\begin{edfigure}[.5em][1em]    	%
\includegraphics[scale=1]{%
gesamttex/edit_VIII-4/images/LH_35_10_16_001-002_d3_001v.pdf}
\edcaption[0.25em]{\lbrack\textit{Fig.~3}\rbrack}
\end{edfigure}
\pstart
(4) Aliter per mera mechanica,\indexS{mechanica} nullo rei physicae occultae\indexS{res physica occulta} (qualis magnetica vis\indexS{vis magnetica} est) interventu idem assequemur. 
\edtext{Esto recta rigida \textit{CR}, mobilis circa punctum \textit{N}.}{\lemma{Esto}\Bfootnote{%
\textit{(1)}~ut supr
\textit{(2)}~recta
\bars{rigida~\textit{erg.}}
\textit{CR}, %
\textit{(a)}~et aliquod punctum \textit{D}, ex quo in rectam agatur normalis \textit{DN}. Po
\textit{(b)}~mobilis circa punctum \textit{N}.~\textit{L}}} 
Extremitati
\edtext{\textit{C}}{\lemma{}\Bfootnote{\textit{C} \textit{erg.~L}}} 
alligetur filum\indexS{filum} \textit{CDEF} 
\edtext{transiens per \textit{D} et super trochlea\indexS{trochlea} \textit{E},}{\lemma{transiens}\Bfootnote{%
\textit{(1)}~ super trochleam \textit{E}, %
\textit{(2)}~per \textit{D} et super %
\textit{(a)}~trochleam %
\textit{(b)}~trochlea \textit{E},~\textit{L}}} 
in extremo appensum 
\edtext{habens pondus\indexS{pondus filo appensum} \textit{F}, quod incumbat plano inclinato\indexS{planum inclinatum} \textit{G}. Similiterque}{\lemma{habens}\Bfootnote{%
\textit{(1)}~filum %
\textit{(2)}~pondus \textit{F} %
\textit{(a)}~incumbens %
\textit{(b)}~, quod incumbat plano inclinato \textit{G}.
\textit{(aa)}~Ac %
\textit{(bb)}~Similiterque~\textit{L}}} 
extremitati \textit{R}, allegetur filum\indexS{filum} \textit{RDHI}, transiens 
\edtext{per \textit{D}, vel}{%
\lemma{per \textit{D},}%
\Bfootnote{%
\textit{(1)}~et su
\textit{(2)}~vel~\textit{L}}}
proxime ad \textit{D}, et 
\edtext{super trochlea\indexS{trochlea}}{\lemma{super}\Bfootnote{%
\textit{(1)}~trochleam %
\textit{(2)}~trochlea~\textit{L}}} 
\textit{H}, in extremo appensum habens pondus\indexS{pondus filo appensum} \textit{I}, quod incumbat plano inclinato\indexS{planum inclinatum} \textit{K}. His positis patet \textit{D} respondere centro 
\edtext{gravium Neostatico,\indexS{centrum Neostaticum gravium} et}{\lemma{gravium}\Bfootnote{%
\textit{(1)}~,~et po %
\textit{(2)}~Geostatico %
\textit{(3)}~Neostatico, et~\textit{L}}} 
pondera \textit{F} et \textit{I}, respondere  
\edtext{ponderibus Neostatice appensis\indexS{pondera Neostatice appensa} in \textit{C} et \textit{R}, et perinde agere, ac si Neostatice essent pondera affixa\indexS{pondera Neostatice affixa} in \textit{C} et \textit{R}}{\lemma{ponderibus}\Bfootnote{%
\bars{\textit{(1)}~Neosta %
\textit{(2)}~Geostatice %
\textit{(3)}~Neostatice~\textit{erg.}}
appensis in \textit{C} et \textit{R},
\textit{(1)}~tendentibus versus \textit{D} %
\textit{(2)}~tractis versus \textit{D} %
\textit{(3)}~et perinde agere %
\bars{in \textit{C} et \textit{R} \textit{gestr.}}, ac si %
\textit{(a)}~essent %
\textit{(b)}~Neostatice essent pondera %
\textit{(aa)}~appensa in %
\textit{(bb)}~affixa in %
\textit{(aaa)}~\textit{R} et \textit{C} %
\textit{(bbb)}~\textit{C} et \textit{R}~\textit{L}}}
tendentiaque versus \textit{D}. Gravitationes autem 
\edtext{qua\lbrack s\rbrack}{%
\lemma{}%
\Bfootnote{%
quam 
\textit{L ändert Hrsg.}}}
habebunt 
\edtext{pondera \textit{F} et \textit{I},}{\lemma{pondera}\Bfootnote{%
\textit{(1)}~\textit{I} et \textit{K} %
\textit{(2)}~\textit{F} et \textit{I},~\textit{L}}} 
in planis inclinatis\indexS{gravitatio in plano inclinato} \textit{G} et \textit{K}, respondebunt solicitationibus ponderum 
\edtext{Neostaticis\indexS{solicitatio Neostatica ponderis} versus}{\lemma{Neostaticis}\Bfootnote{%
\textit{(1)}~ad %
\textit{(2)}~versus centrum \textit{D} %
\textit{(a)}~seu %
\textit{(b)}~et poterunt decrescere %
\textit{(aa)}~esse %
\textit{(aaa)}~tanto majores, 
\textit{(bbb)}~ma %
\textit{(bb)}~decrescere~\textit{L}}} 
crescentibus distantiis, ita ut crescente \textit{CD}, decrescat declivitas ipsius plani,\indexS{declivitas plani} \textit{G}, vel decrescente \textit{RD}, crescat declivitas ipsius plani\indexS{declivitas plani} \textit{K}. 
Itaque in Neostaticis etiam poterunt solicitationes crescere decrescentibus distantiis,\indexS{solicitatio decrescente distantia crescens} et decrescere crescentibus.\indexS{solicitatio crescente distantia decrescens} Et 
\edtext{quod dici}{\lemma{quod}\Bfootnote{%
\textit{(1)}~de %
\textit{(2)}~dici~\textit{L}}} 
poterit in nostra Machinatione Staticae vulgaris\indexS{Statica vulgaris} hic repraesentata, locum etiam habebit in Hypothesi Neostaticae\indexS{Neostatica} qualicunque. 
\pend 
\pstart
\noindent
\href{\Scan{LH_35_10_16_002r}}%
{\lbrack2~r\textsuperscript{o}\rbrack}\  
\pend
\nobreak
\begin{edfigure}[1.0em]%
\includegraphics[scale=1]{%
gesamttex/edit_VIII-4/images/LH_35_10_16_001-002_d4_002r.pdf}
\edcaption{%
\edtext{\lbrack\textit{Fig.~4}\rbrack}{\lemma{%
\lbrack\textit{Fig.~4}\rbrack}\Cfootnote[nonum]{Vorlage ist die Figur \cite{00541}a.\,a.\,O., S.\;67.
Ein gestrichener Entwurf zum Diagramm wird nicht wiedergegeben.}}}
\end{edfigure}
\pstart
(5) Sed operae pretium 
\edtext{erit aestimationem ipsam nonnihil}{\lemma{erit}\Bfootnote{%
\textit{(1)}~rem nonnihil 
\textit{(a)}~ad %
\textit{(b)}~calculo %
\textit{(2)}~aestimationem ipsam nonnihil~\textit{L}}} 
attentare. Ex puncto \textit{N} in rectas \textit{CD}, \textit{RD}, agantur normales \textit{NL}, \textit{NX}. 
\edlabel{35_10_16_001-002_5a}%
\edtext{}{%
\xxref%
{35_10_16_001-002_5a}{35_10_16_001-002_5b}%
\lemma{Constat \lbrack...\rbrack\ perpendiculares}%
\applabel{FnC:35_10_16_003-004_alibi-ostensis}%
\Cfootnote{%
Siehe bspw. die \textit{Solutio problematis R.\,P.\;Vannii S.\,J.} (April bis Juni 1690), N.\;\ref{41823} in diesem Band, bes. \refpassage{35_10_13_010-011_17a}{35_10_13_010-011_17b}. 
Der in der Folge für diesen Satz erbrachte Beweis 
gilt nur unter der Annahme, dass die Dreiecke \textit{NCD} und \textit{NRD} ähnlich sind, d.\,h., dass der Winkel \textit{D} recht ist; eine Annahme, die zwar nicht an dieser Stelle, aber in der entsprechenden Passage der Vorlage ausdrücklich gemacht wird (siehe \refnote{FnB:LH_35_10_16_004_angulo-recto} von  N.\;\ref{60124}).}}%
Constat (ex alibi
\edtext{ostensis,) fore aequilibrium\indexS{aequilibrium} in \textit{N}, si sint momenta reciproce ut perpendiculares;\edlabel{35_10_16_001-002_5b} seu si sit}{\lemma{ostensis,}\Bfootnote{%
\textit{(1)}~quae %
\textit{(a)}~nec %
\textit{(b)}~etiam %
\textit{(c)}~\bars{ex \textit{streicht Hrsg.}} prop. 6. libri 2. \textit{Neostaticae} %
\textit{(aa)}~abludunt) fore aequilib %
\textit{(bb)}~duci possent) %
\textit{(2)}~)~fore aequilibrium in \textit{N},
\textit{(1)}~si sit momentum in \textit{C} %
\textit{(2)}~si \lbrack...\rbrack\ seu \bars{si~\textit{erg.}} sit~\textit{L}}} 
momentum in \textit{C} ad momentum in \textit{R}, ut \textit{NX} ad \textit{NL}. 
\edtext{%
\href{\Scan{LH_35_10_16_001v}}%
{\lbrack1~v\textsuperscript{o}\rbrack}\ 
Idque etiam consentit in effectu cum 
\edtext{propositione sexta libri secundi \textit{Neostaticae}.}{\lemma{propositione sexta libri secundi \textit{Neostaticae}}\Cfootnote{\indexN{\textls{Saccheri}, Girolamo 1667\textendash 1733}\textsc{Saccheri}, \cite{00541}\textit{Neo-statica}, S.\;59\,f.}}
Est enim \textit{NC} ad \textit{NR} ut triang. \textit{NCD} ad triang. \textit{NRD}, id est ut \textit{NL} in \textit{CD} ad \textit{NX} in \textit{RD}. Ergo fit \textit{NC} in \textit{RD} ad \textit{NR} in \textit{CD}, ut \textit{NL} ad \textit{NX}.}{\lemma{}\Bfootnote{%
\href{\Scan{LH_35_10_16_001v}}%
{\lbrack1~v\textsuperscript{o}\rbrack}\ 
Idque \lbrack...\rbrack\  \textit{Neostaticae}. %
\textit{(1)}~Est enim %
\textit{(a)}~rectang. %
\textit{(b)}~triang. \textit{DCN} ad triang. %
\textit{(c)}~\textit{CN} ad %
\textit{(d)}~\textit{NC} ad \textit{NR} %
\textit{(2)}~Est enim \lbrack...\rbrack\ ad \textit{NX}. \textit{erg.~L}}} 
\href{\Scan{LH_35_10_16_002r}}%
{\lbrack2~r\textsuperscript{o}\rbrack}\  
Momentum\indexS{momentum} autem 
\edtext{mihi hic}{\lemma{}\Bfootnote{mihi hic \textit{erg.~L}}} 
significat factum ex pondere in solicitationem,\indexS{solicitatio}\indexS{factum ex pondere in solicitationem} seu 
\edtext{momenta\indexS{momentum} mihi hic sunt}{\lemma{momenta}\Bfootnote{%
\textit{(1)}~sunt %
\textit{(2)}~mihi hic sunt~\textit{L}}} 
in ratione composita ponderum et 
\edtext{solicitationum, vel quod idem est,}{\lemma{solicitationum,}\Bfootnote{%
\textit{(1)}~seu %
\textit{(2)}~vel quod idem est,~\textit{L}}} 
momentum in \textit{C}, est ad momentum in \textit{R},
\edtext{in ratione composita ponderis in \textit{C} ad pondus in \textit{R}, et solicitationis in \textit{C} ad solicitationem\indexS{solicitatio}
\edlabel{35_10_16_001-002_1a}%
\edtext{}{%
\xxref%
{35_10_16_001-002_1a}{35_10_16_001-002_1b}%
\lemma{in \textit{R}.}\Bfootnote{%
\textit{(1)}~Ponamus jam puncta \textit{R}, \textit{N}, \textit{D} esse data, item pondera in \textit{C} et in \textit{R} ponenda \bars{esse data \textit{erg.}}; quaeritur an punctum \textit{C} obtineri possit, posito solicitationes paracentricas esse minores in distantiis \bars{a centro~\textit{erg.}} majoribus. %
\textit{(2)}~(6) %
\textit{(a)}~Hic ante omnia %
\textit{(b)}~Praeterea~\textit{L}}}%
in \textit{R}.}{\lemma{}\Afootnote{%
\textit{Am Rand, gestrichen}:\enskip 
{\footnotesize
$CN:NR=CD\cdot NL:RD\cdot NX$. Ergo\textsuperscript{\lbrack a\rbrack} $CN\cdot RD:NR\cdot CD=NL:NX$.
\\[1.5ex] %
\textsuperscript{\lbrack a\rbrack}
\textit{(1)}~$CN\cdot NX:N$ %
\textit{(2)}~$CN\cdot RD$~\textit{L}}}}
\pend 
\begin{edfigure}%
\hspace*{-1em}%
\includegraphics[scale=1]{%
gesamttex/edit_VIII-4/images/LH_35_10_16_001-002_d5_002r.pdf}
\edcaption{\edtext{\lbrack\textit{Fig.~5}\rbrack}{\lemma{\lbrack\textit{Fig.~5}\rbrack}\Cfootnote[nonum]{Vorlage ist die Figur \cite{00541}a.\,a.\,O., S.\;67.}}}
\end{edfigure}
\pstart
(6) Praeterea\edlabel{35_10_16_001-002_1b} observo, 
\edtext{quaecunque affirmatur hypothesis quantitatis solicitationum\indexS{quantitas solicitationis}\lbrack,\rbrack\ si angulus \textit{DNR} sit rectus\lbrack,\rbrack}{\lemma{}%
\Bfootnote{quaecunque \lbrack...\rbrack\  rectus \textit{erg.~L}}} 
dato momento\indexS{momentum} in \textit{R}, posse quidem 
\edtext{momentum in \textit{C},}{\lemma{momentum}\Bfootnote{%
\textit{(1)}~\textit{C} %
\textit{(2)}~in \textit{C},~\textit{L}}} 
concipi 
\edtext{quantumvis eo majus momento in \textit{R} quia}{\lemma{quantumvis}\Bfootnote{%
\textit{(1)}~minus %
\textit{(2)}~eo majus
\bars{momento
\textit{(1)}~in \textit{R} %
\textit{(2)}~in \textit{R} \textit{erg.}} quia~\textit{L}}} 
etiam ipsa
\edtext{\textit{NX} data,}{\lemma{\textit{NX}}\Bfootnote{%
\textit{(1)}~d %
\textit{(2)}~potest esse %
\textit{(3)}~data,~\textit{L}}} 
potest esse in quantavis ratione major 
\edtext{ipsa \textit{NL}, cum}{\lemma{ipsa \textit{NL},}\Bfootnote{%
\textit{(1)}~quia %
\textit{(2)}~cum~\textit{L}}} 
\textit{C} tam prope 
\edtext{admoveri possit}{\lemma{admoveri}\Bfootnote{%
\textit{(1)}~potest %
\textit{(2)}~possit~\textit{L}}} 
ad \textit{N}, ut \textit{NL} fiat quantumvis 
\edtext{parva. Sed observo vicissim momentum in \textit{C} non posse concipi in quantavis ratione minus}{\lemma{parva.}\Bfootnote{%
\textit{(1)}~Sed momentum in \textit{C} non potest %
\textit{(2)}~Sed \lbrack...\rbrack\ concipi 
\textit{(a)}~quantumvis minus %
\textit{(b)}~in quantavis ratione minus~\textit{L}}} 
momento in \textit{R}, quia ipsa \textit{NL} non potest assumi quantumvis major data \textit{NX}, 
\edtext{cum \textit{NL} necessario sit minor ipsa \textit{ND}\lbrack,\rbrack}{\lemma{cum}\Bfootnote{%
\textit{(1)}~non possit %
\textit{(a)}~esse major quam \textit{ND} %
\textit{(b)}~pervenire \textit{NL} ad aequalitatem ipsius \textit{ND} %
\textit{(2)}~\textit{NL} necessario sit minor %
\textit{(a)}~ipsa \textit{ND} perpendiculari ex \textit{D} in \textit{CR} ducta %
\textit{(b)}~ipsa \textit{ND},~\textit{L}}} 
quantumvis \textit{C} removeatur ab \textit{N}. Itaque 
\edtext{hic}{\lemma{}\Bfootnote{hic \textit{erg.~L}}} 
assumtio rationis inter momenta,\indexS{momentum} vel datis ponderibus, assumtio rationis inter 
\edlabel{35_10_16_001-002_3a}%
\edtext{}{\xxref{35_10_16_001-002_3a}{35_10_16_001-002_3b}\lemma{solicitationes,}%
\Bfootnote{%
\textit{(1)}~aliquas %
\textit{(2)}~aliquos limites %
\textit{(a)}~reci %
\textit{(b)}~recipit %
\textit{(c)}~postulat.
\lbrack/\rbrack\ 
\textit{(aa)}~Quaeritur jam %
\textit{(aaa)}~datis \textit{R} et %
\textit{(bbb)}~positis utcunque punctis, \textit{R}, \textit{N}, \textit{D}, angulo \textit{RND} recto, %
\textit{(aaaa)}~datis etiam ponderibus %
\textit{(bbbb)}~ponderibusque in \textit{C} et in \textit{R} ponendis, positoque %
\textit{(bb)}~(7) Ponamus jam~\textit{L}}}%
solicitationes,\indexS{solicitatio} aliquos limites postulat.
\pend 
\pstart 
(7) Ponamus jam\edlabel{35_10_16_001-002_3b} 
solicitationes\indexS{solicitatio in distantia majori minor} esse minores in distantiis 
\edtext{majoribus; quaeritur regula aequilibrii.\indexS{regula aequilibrii} Redeamus ad 
\edtext{figuram paragraphi 5.}{%
\lemma{figuram paragraphi 5}%
\Cfootnote{Vgl. das Diagramm \lbrack\textit{Fig.~4}\rbrack.}}
Pondus}{\lemma{majoribus;}\Bfootnote{%
\textit{(1)}~vel contra, %
\textit{(2)}~quaeritur 
\textit{(a)}~quomodo %
\textit{(aa)}~desi %
\textit{(bb)}~inveniri possit %
\textit{(aaa)}~punctum %
\textit{(aaaa)}~\textit{D} %
\textit{(bbbb)}~\textit{C} %
\textit{(cc)}~regula %
\textit{(dd)}~regula aequilibrii 
\bars{posito 
\textit{(1)}~centrum %
\textit{(2)}~\bars{punctum \textit{streicht Hrsg.}} aequilibrii esse \textit{N} \textit{erg.~u.~gestr.}}.
\bars{Redeamus ad figuram paragraphi 5.~\textit{erg.}}
Pondus~\textit{L}}} 
in \textit{C}, vocetur \textit{c}, solicitatio in \textit{C}, vocetur \textit{k}\lbrack,\rbrack\ pondus in \textit{R} vocetur \textit{r}, et solicitatio in \textit{R}, 
\edtext{vocetur $\rho $. Erunt momenta ut \textit{ck}, $r\rho $, et}{\lemma{vocetur $\rho $.}\Bfootnote{%
\textit{(1)}~$NR\cdot CD:RD$ %
\textit{(2)}~Erunt momenta \bars{ut \textit{erg.}} \textit{ck}, $r\rho $, et~\textit{L}}} 
fiet 
\edtext{\textit{ck} ad $r\rho $ ut}{\lemma{\textit{ck} ad}\Bfootnote{%
\textit{(1)}~\textit{rR} %
\textit{(2)}~$r\rho $ %
\textit{(a)}~$=$ %
\textit{(b)}~ut~\textit{L}}} 
\textit{NX} ad \textit{NL} (ex principiis\indexS{principia mechanica} 
\edtext{mechanicis). Sed \textit{NX} est ad \textit{NL} ut \textit{NR} in \textit{CD} ad \textit{NC} in \textit{RD}}{\lemma{mechanicis)}\Bfootnote{%
\textit{(1)}~$=$ %
\textit{(2)}~seu ut \textit{CN} \bars{in \textit{erg.}} $RD:NR$ \bars{in \textit{erg.}} \textit{CD}  %
\textit{(3)}~. Sed \textit{NX} est ad \textit{NL} ut \textit{NR} \bars{in \textit{erg.}} \textit{CD} ad \textit{NC} \bars{in \textit{erg.}} \textit{RD}~\textit{L}}} 
(ex Geometria.\indexS{Geometria}) 
\edtext{Itaque}{\lemma{}\Bfootnote{Itaque \textit{erg.~L}}} 
fiet \textit{ck} 
\edtext{ad $r\rho $ ut \textit{NR} in \textit{CD} ad \textit{NC} in \textit{RD}}{\lemma{ad $r\rho $}\Bfootnote{%
\textit{(1)}~$=$ %
\textit{(2)}~ut %
\textit{(a)}~\textit{CN} \bars{in \textit{erg.}} $RD:RN$ in \textit{CD} %
\textit{(b)}~\textit{NR} \bars{in \textit{erg.}} \textit{CD} ad \textit{NC} \bars{in \textit{erg.}} \textit{RD}~\textit{L}}} 
seu fiet
\edlabel{35_10_16_001-002_4a}%
\edtext{}{\xxref{35_10_16_001-002_4a}{35_10_16_001-002_4b}\lemma{\textit{k} in \textit{RD} ad}%
\Bfootnote{%
\textit{(1)}~$\rho RD$ %
\textit{(2)}~$\rho $ in \textit{CD} %
\textit{(a)}~$=$ %
\textit{(b)}~ut %
\textit{(aa)}~\textit{CN} %
\textit{(bb)}~\textit{r} in $CN:$ %
\textit{(cc)}~\textit{r} in \textit{NR} ad \textit{c} in %
\textit{(aaa)}~\textit{RN} %
\textit{(bbb)}~\textit{NC}.
\textit{(aaaa)}~Si ergo crescente \textit{CD} decrescit \textit{k} et crescente \textit{RD} decrescit $\rho $ %
\textit{(bbbb)}~Si jam ponamus %
\textit{(aaaaa)}~\textit{k} %
\textit{(bbbbb)}~solicitationes %
\textit{(aaaaa-a)}~crescere ut %
\textit{(bbbbb-b)}~esse %
\textit{(cccc)}~V %
\textit{(dddd)}~\href{\Scan{LH_35_10_16_002v}}%
{\lbrack2~v\textsuperscript{o}\rbrack}\ (8) Veniamus
\textit{(aaaaa)}~jam %
\textit{(bbbbb)}~nunc~\textit{L}}}%
\textit{k} in \textit{RD} ad $\rho $ in \textit{CD} ut \textit{r} in \textit{NR} ad \textit{c} in \textit{NC}.
\pend
\pstart
\href{\Scan{LH_35_10_16_002v}}%
{\lbrack2~v\textsuperscript{o}\rbrack}\ 
(8) Veniamus nunc\edlabel{35_10_16_001-002_4b}
ad hypothesin 
\edtext{aliquam}{%
\lemma{}%
\Bfootnote{%
aliquam 
\textit{erg.~L}}} 
determinatam, et ponamus solicitationes esse reciproce ut distantias,\indexS{solicitationes reciprocae distantiis} 
\edtext{seu \textit{k} ad $\rho $ ut}{\lemma{seu}\Bfootnote{%
\textit{(1)}~$k\rho $ %
\textit{(2)}~\textit{k} ad $\rho $ %
\textit{(a)}~= %
\textit{(b)}~ut~\textit{L}}} 
\textit{RD} 
\edtext{ad \textit{CD}. Fiet \textit{k} in \textit{RD} ad $\rho $ in \textit{CD} ut quadr. \textit{RD} ad quadr. \textit{CD}, quod erit proinde ut \textit{r} in \textit{NR} ad \textit{c} in \textit{NC} per paragraphum praecedentem. Et erunt etiam pondera (\textit{c} et \textit{r}) in ratione composita}{\lemma{ad \textit{CD}}\Bfootnote{%
\textit{(1)}~,~fiet ratio \textit{k} in \textit{RD} ad $\rho $ in %
\textit{(a)}~\textit{RD} %
\textit{(b)}~\textit{CD} ratio aequalitatis, adeoque ratio %
\textit{(aa)}~\textit{r} in \textit{CN}, ad \textit{c} in \textit{RN} %
\textit{(bb)}~\textit{r} in \textit{NR} ad \textit{c} in \textit{NC} erit %
\textit{(aaa)}~ratio %
\textit{(bbb)}~etiam ratio aequalitatis, seu %
\textit{(aaaa)}~\textit{r} in %
\textit{(aaaaa)}~\textit{C} %
\textit{(bbbbb)}~\textit{NR} erit %
\textit{(bbbb)}~erit \textit{r} in $NR=$ %
\textit{(aaaaa)}~\textit{c} in \textit{RN} %
\textit{(bbbbb)}~\textit{c} in \textit{CN} %
\textit{(aaaaa-a)}~seu ut in communi statica, erunt ponde %
\textit{(bbbbb-b)}~seu erunt 
\textit{(ccccc-c)}~. Et erit pondus in \textit{C} ad pondus in \textit{R}, ut \textit{CN} ad \textit{RN} seu erit pondus in \textit{C} ad pondus in \textit{R}, ut %
\textit{(aaaaa-aa)}~\textit{RC} %
\textit{(bbbbb-bb)}~\textit{NR} ad \textit{NC} prorsus ut in statica vulgari %
\textit{(aaaaa-aaa)}~. Ibi enim %
\textit{(bbbbb-bbb)}~quae exhibet revera casum specialem hypotheseos nostrae. Cum enim in statica vulgari distantiae ponantur infinitae aequales, %
\textit{(aaaaa-aaaa)}~fient %
\textit{(bbbbb-bbbb)}~itaque si solicitationes sint reciproce ut distantiae fient etiam aequales. %
\textit{(2)}~.~Fiet \textit{k} \lbrack...\rbrack\ quadr. \textit{CD}, %
\textit{(a)}~quod erit $=$ %
\textit{(b)}~quod erit \lbrack...\rbrack\ ad \textit{c} 
\textit{(aa)}~in
\textit{(bb)}~in \textit{NC} per paragraphum praecedentem. Et
\textit{(aaa)}~erit etiam pondus in \textit{C} ad pondus in \textit{R} in %
\textit{(bbb)}~erunt etiam pondera %
\textit{(aaaa)}~(\textit{r} et \textit{c}) %
\textit{(bbbb)}~(\textit{r} et \textit{c}) %
\textit{(cccc)}~(\textit{c} et \textit{r}) in ratione composita~\textit{L}}} 
ex brachiorum reciproca simplice,\indexS{ratio brachiorum reciproca simplex} 
et distantiarum directa duplicata,\indexS{ratio distantiarum duplicata} 
seu erit pondus in \textit{C} ad pondus in \textit{R}, ut \textit{NR} in quadr. \textit{CD} ad \textit{NC} in quadr. \textit{RD}. 
Hinc quia hypothesis staticae vulgaris\indexS{statica vulgaris} haberi potest pro casu 
\edtext{speciali}{\lemma{}\Bfootnote{speciali \textit{erg.~L}}} 
hypotheseos 
\edtext{hujus praesentis,}{\lemma{hujus}\Bfootnote{%
\textit{(1)}~nostrae %
\textit{(2)}~praesentis,~\textit{L}}} 
quoniam tunc distantiae infinitae\indexS{distantia infinita} habentur pro aequalibus, adeoque et 
\edtext{solicitationes (si}{\lemma{solicitationes (}\Bfootnote{%
\textit{(1)}~cum %
\textit{(2)}~si~\textit{L}}} 
sint distantiis 
\edtext{reciprocae) fiunt aequales; hinc patet}{\lemma{reciprocae}\Bfootnote{%
\textit{(1)}~fient aequales ) %
\textit{(2)}~) fiunt aequales; hinc patet~\textit{L}}} 
rationem quadr. \textit{CD} ad quadr. \textit{RD} fore 
\edtext{tunc}{\lemma{}\Bfootnote{tunc \textit{erg.~L}}} 
aequalitatis ac proinde 
\edtext{tunc}{\lemma{}\Bfootnote{tunc \textit{erg.~L}}} 
fore pondus in \textit{C} ad pondus in \textit{R} ut \textit{NR} 
\edtext{ad \textit{NC}, uti esse debebat.}{\lemma{ad \textit{NC},}\Bfootnote{%
\textit{(1)}~quod %
\textit{(2)}~quemadmodum 
\textit{(a)}~constat. %
\textit{(b)}~et aliunde constat. %
\textit{(3)}~uti esse debebat.~\textit{L}}} 
\pend 
\pstart 
(9) Quod 
\edtext{si assumamus}{\lemma{si}\Bfootnote{%
\textit{(1)}~ponamus %
\textit{(2)}~assumamus~\textit{L}}} 
aliam Hypothesin, et ut arbitror veriorem, solicitationes\indexS{solicitationes reciprocae quadratis distantiarum} esse reciproce ut quadrata 
\edtext{distantiarum,\indexS{quadratum distantiae} seu \textit{k} ad $\rho $ esse ut qu. \textit{RD} ad qu. \textit{CD}, fient (ex paragrapho~7.) pondera}{\lemma{distantiarum,}\Bfootnote{%
\textit{(1)}~fiet %
\textit{(2)}~seu \textit{k} ad $\rho $ %
\textit{(a)}~$=$ %
\textit{(b)}~esse ut qu. \textit{RD} ad qu. \textit{CD}, %
\textit{(aa)}~et %
\textit{(bb)}~fiet %
\textit{(cc)}~fient (ex paragrapho
\textit{(aaa)}~penultima\lbrack\textit{!}\rbrack\ %
\textit{(bbb)}~antepenultima\lbrack\textit{!}\rbrack). %
\textit{(ccc)}~7.) %
\textit{(aaaa)}~\textit{k} in $RD:\rho \cdot CD$ %
\textit{(bbbb)}~fient %
\textit{(cccc)}~pondera~\textit{L}}} 
in ratione 
\edtext{composita ex brachiorum}{\lemma{composita ex}\Bfootnote{%
\textit{(1)}~ponderum %
\textit{(2)}~brachiorum~\textit{L}}} 
reciproca simplice\indexS{ratio brachiorum reciproca simplex} et distantiarum directa triplicata,\indexS{ratio distantiarum triplicata} 
seu erit pondus in \textit{C} ad pondus in \textit{R} ut \textit{NR} in cub. \textit{CD} 
ad \textit{NC} in cub. \textit{RD}. quae Hypothesis, 
etiam in vulgarem\indexS{Hypothesis vulgaris} evanescit, cum distantiae ponuntur 
\edlabel{35_10_16_001-002_2a}%
\edtext{}{\xxref{35_10_16_001-002_2a}{35_10_16_001-002_2b}\lemma{infinitae}\Bfootnote{%
\textit{(1)}~adeoque aequa %
\textit{(2)}~pro aequalibus habendae.
\textit{(a)}~Quae cum ita sint, operae pretium erit accuratius expendi quae allata sunt %
\textit{(aa)}~ad a %
\textit{(bb)}~in \textit{Neostatica} ad astruendum corollarium propositionis 15 libri secundi. %
\textit{(b)}~(10) Quod si~\textit{L}}}
infinitae\indexS{distantia infinita} pro aequalibus habendae.
\pend 
\pstart
(10) Quod si\edlabel{35_10_16_001-002_2b}
distantiis positis inaequalibus nihilominus solicitationes essent aequales, 
\edtext{fierent pondera in ratione composita ex brachiorum reciproca\indexS{ratio brachiorum reciproca}}{%
\lemma{fierent}%
\Bfootnote{%
\textit{(1)}~in %
\textit{(2)}~pondera in ratione composita ex brachiorum %
\textit{(a)}~directa, %
\textit{(b)}~reciproca~\textit{L}}}
et distantiarum directa.\indexS{ratio distantiarum directa} 
Denique, si solicitationes essent ut 
\edtext{distantiae secundum hypothesin nostri autoris pondera forent}{%
\lemma{distantiae}%
\Bfootnote{%
\bars{secundum hypothesin nostri autoris \textit{erg.}}
\textit{(1)}~fierent pon
\textit{(2)}~pondera forent~\textit{L}}}
reciproce ut brachia, nec regula aequilibrii\indexS{regula aequilibrii} differret a regula staticae vulgaris,\indexS{statica vulgaris}
\edtext{quod etiam Autor clarissimus bene observavit. Sed}{%
\lemma{quod}%
\Bfootnote{%
\textit{(1)}~etiam Cl. Autor \textit{Neostaticae} observavit. Sed %
\textit{(2)}~etiam %
\textit{(a)}~vir %
\textit{(b)}~Autor clarissimus bene observavit. Sed~\textit{L}}}
ex caeteris hactenus comprobatis, apparet accuratius expendenda esse, quae allate sunt in \textit{Neostatica} ad astruendum 
\edtext{Corollarium propositionis 15 libri secundi.}{\lemma{Corollarium propositionis 15 libri secundi}\Cfootnote{\indexN{\textls{Saccheri}, Girolamo 1667\textendash1733}\textsc{Saccheri}, \cite{00541}\textit{Neo-statica}, S.\;69.}} 
\pend 
